\section{Compra y producción propia de datos: ¿comprar datos o crearlos?}

En la industria, las opciones habituales son comprar datos, recolectar datos y crear datos. Comprar datos es rápido, pero lleva a la homogeneidad; construir datos propios crea barreras profundas, pero es costoso; los datos sintéticos o de simulación son escalables, pero conllevan un riesgo de falta de realismo. La elección cambia según la etapa: en la validación temprana se puede comprar o sintetizar, pero al llegar a las capacidades centrales es indispensable construir un ciclo cerrado propio.

\subsection{Comparación de las tres vías}

\noindent\begin{tabularx}{\textwidth}{p{0.20\textwidth}p{0.34\textwidth}X}
\toprule
Vía & Etapa adecuada & Riesgos \\
\midrule
Comprar datos & Arranque rápido, complemento de capacidades generales & Homogeneidad, licencias poco claras, rendimiento marginal bajo. \\
Recolectar datos & Desarrollo de capacidades en tareas reales & Costo alto, plazos largos, privacidad y seguridad complejas. \\
Crear datos & Escenarios de cola larga, peligrosos o escasos & Desplazamiento de la distribución, evaluaciones distorsionadas, sobreajuste a la simulación. \\
\bottomrule
\end{tabularx}

\begin{importantbox}{Las barreras de largo plazo provienen de un ciclo cerrado propio}
Los datos que se pueden comprar difícilmente se convierten en una ventaja defensiva de largo plazo. La verdadera barrera proviene de la combinación de tareas propias, retroalimentación propia, evaluación propia y Recipe propia.
\end{importantbox}

\subsection{Resumen de la sección}

La estrategia de datos no consiste en elegir entre comprar o crear. A corto plazo se puede comprar y sintetizar; a largo plazo es indispensable formar un ciclo cerrado de datos y un sistema de evaluación propios.
